\documentclass[11pt,a4paper]{article}
\usepackage[T1]{fontenc}
\usepackage[utf8]{inputenc}
\usepackage{lmodern,amsmath,amssymb}
\usepackage[margin=25mm]{geometry}
\usepackage{longtable,booktabs,array,calc}
\usepackage[hidelinks]{hyperref}
\hypersetup{pdftitle={The zero-spacing limit for any plaquette action and any dimension},pdfauthor={navstokgap research working notes}}
\newcommand{\tightlist}{\setlength{\itemsep}{0pt}\setlength{\parskip}{0pt}}
\setlength{\emergencystretch}{3em}
\setcounter{secnumdepth}{0}
\begin{document}
\hypertarget{the-zero-spacing-limit-for-any-plaquette-action-and-any-dimension}{%
\section{The zero-spacing limit for any plaquette action and any
dimension}\label{the-zero-spacing-limit-for-any-plaquette-action-and-any-dimension}}

\textbf{Result, 2026-09-27.} The
\href{series-parallel-gauge-refinement.md}{series/parallel note} worked
with the heat-kernel action. Here the question is how the \(a\to0\)
limit depends on the choice of plaquette action, and on the dimension.

\begin{enumerate}
\def\labelenumi{\arabic{enumi}.}
\tightlist
\item
  \textbf{Two semigroups.} A series move convolves plaquette weights, so
  it multiplies their character coefficients; the heat kernel is the
  family closed under it. A parallel move multiplies weights pointwise,
  so it adds their logarithms; exponential families \(e^{-\beta V}\)
  (Wilson, Villain-type) are closed under it. In the small-field regime
  the two agree: for \(U(1)\) the product of two heat kernels is exactly
  a heat kernel at the parallel heat time \(st/(s+t)\) plus vortex terms
  of relative size \(e^{-2\pi^2/(s+t)}\) (Proposition 1).
\item
  \textbf{Two dimensions, any action} (Theorem 2). The possible
  continuum limits of refining a two-dimensional lattice gauge theory
  with any symmetric class-function plaquette weight are exactly the
  conjugation-invariant Lévy exponents
  \(\psi(R)=\frac{\sigma^2}2C_2(R)+\int(1-{\rm Re}\,\chi_R/d_R)\,d\nu\),
  with \(\psi(R)=\lim a^{-2}(1-\hat c_R(a))\). Yang--Mills is the
  Gaussian member, \(\nu=0\), selected by a Lindeberg condition on the
  action. A weight with an atom at a centre element produces a
  centre-vortex gas, and the circle spectrum becomes
  \(E_R-E_0=\hbar cL\,\psi(R)\).
\item
  \textbf{Any dimension: an error budget per physical volume.} One
  refinement step leaves three kinds of error. The perturbative part is
  of relative order \(t_n\), the plaquette heat time at step \(n\). The
  large-field part has density \(a_n^{-D}e^{-c/t_n}\) per physical
  volume, with \(c\) the minimal large-field action of the chosen
  lattice action in units of the inverse heat time. The jump part is the
  action's Lévy content. With \(t\propto\lambda_Da^{4-D}\):

  \begin{itemize}
  \tightlist
  \item
    \(D<4\): all three are summable at fixed \(\lambda_D\) once the jump
    part vanishes, which is why the lattice gap of compact \(U(1)\)
    disappears at fixed coupling (Gross, Theorem 5 of the
    series/parallel note);
  \item
    \(D=4\): the perturbative part resums into the running coupling, and
    the large-field part decays only as a power of \(a\),
    \((a\Lambda)^{2b_0c}\), so it dies iff \(2b_0c>4\), that is
    \(c>96\pi^2/(11N)\) for pure \(SU(N)\): \(c>43.1\) for \(SU(2)\),
    \(c>28.7\) for \(SU(3)\). Instantons (\(c=8\pi^2\approx79\)) pass
    for every \(N\), and the condition on lattice dislocations is the
    kind studied by Pugh and Teper;
  \item
    \(D>4\): \(t\) grows as \(a\to0\) at fixed \(\lambda_D\), so no
    small-field regime exists at short distances.
  \end{itemize}
\end{enumerate}

So, in the ultraviolet, the limit depends on the lattice action only
through its second moment (the coupling, and in \(D=4\) the ratio of
\(\Lambda\) parameters), its minimal large-field action and its jump
content. A gap, being an infrared property, is then the same for every
action that meets the three conditions, up to the \(\Lambda\) ratio. The
ingredients are established (character expansions, Poisson summation,
Hunt's classification of convolution semigroups, the one-loop running);
the organization by the two semigroups and the error budget is the
contribution, with no novelty claimed for the parts.

\hypertarget{series-and-parallel-act-on-different-coordinates-of-the-weight}{%
\subsection{1. Series and parallel act on different coordinates of the
weight}\label{series-and-parallel-act-on-different-coordinates-of-the-weight}}

Let \(w\) be a class-function probability density on a compact connected
Lie group \(G\) (Haar measure normalized), symmetric,
\(w(g^{-1})=w(g)\), with character coefficients
\(\hat c_R=d_R^{-1}\int w\,\overline{\chi_R}\,dg\in[-1,1]\), so that
\(w=\sum_Rd_R\hat c_R\chi_R\). Group metric as in the series/parallel
note §1 (\(-\Delta\chi_R=C_2(R)\chi_R\)).

\begin{itemize}
\tightlist
\item
  \textbf{Series.} Integrating out an edge shared by two faces convolves
  their weights, and
  \(\widehat{(w_1*w_2)}_R=\hat c_R(w_1)\hat c_R(w_2)\). The heat kernels
  \(k_t\), \(\hat c_R=e^{-tC_2(R)/2}\), form the one-parameter family
  closed under this product with \(t\) additive.
\item
  \textbf{Parallel.} Two faces with the same holonomy combine by the
  pointwise product \(w_1w_2\), and \(\log(w_1w_2)=\log w_1+\log w_2\).
  The families \(w_\beta\propto e^{-\beta V}\) with fixed \(V\) are
  closed under this operation with \(\beta\) additive; Wilson's
  \(V=1-{\rm Re\,tr}\,U/N\) is one.
\end{itemize}

\textbf{Proposition 1 (parallel product of \(U(1)\) heat kernels).} With
\(k_s(\theta)=\sum_{n\in\mathbb Z}e^{-sn^2/2}e^{in\theta}\) and
\(\tau=st/(s+t)\),

\[k_s(\theta)\,k_t(\theta)=\frac1{\sqrt{2\pi(s+t)}}\sum_{m\in\mathbb Z}
e^{-2\pi^2m^2/(s+t)}\,k_\tau\!\Bigl(\theta+\frac{2\pi ms}{s+t}\Bigr).\]

\emph{Proof.} Poisson summation gives
\(k_s(\theta)=\sum_wG_s(\theta+2\pi w)\) with
\(G_s(x)=(2\pi s)^{-1/2}e^{-x^2/(2s)}\). For Gaussians,
\(G_s(x)G_t(y)=G_{s+t}(x-y)\,G_\tau\bigl((tx+sy)/(s+t)\bigr)\). Put
\(x=\theta+2\pi w\), \(y=\theta+2\pi w'\) and \(m=w'-w\); the sum over
\(w\) at fixed \(m\) is a wrapped Gaussian of variance \(\tau\) centred
at \(-2\pi ms/(s+t)\), which is \(k_\tau\) at the shifted argument by
Poisson summation again. \(\square\)

The \(m=0\) term is the heat kernel at the parallel heat time, the
conductance rule of the series/parallel note §3. The terms \(m\neq0\)
are vortices, of relative weight \(e^{-2\pi^2m^2/(s+t)}\). So the heat
kernel is closed under series moves exactly and under parallel moves up
to \(e^{-c/t}\) corrections; exponential families have the opposite
property. Both reduce to the Gaussian in the small-field regime, where
they agree.

\textbf{A limit of validity for heat-kernel anisotropy.} Equation (1) of
the series/parallel note matches heat times to the continuum action in
the small-field regime. Refining one direction alone, for instance
Euclidean time at fixed spatial lattice, sends the transverse heat times
to infinity. Heat-kernel weights then become flat up to corrections
\(e^{-tC_2/2}\), whereas the Kogut--Susskind limit needs a magnetic
weight \(e^{-a_0V}\) with \(a_0\) the time step. The Hamiltonian limit
is reached within the exponential family. The directional halvings of an
isotropic cycle change heat times only by factors of two, so the
propositions of that note are unaffected.

\hypertarget{two-dimensions-the-continuum-limits-of-every-plaquette-action}{%
\subsection{2. Two dimensions: the continuum limits of every plaquette
action}\label{two-dimensions-the-continuum-limits-of-every-plaquette-action}}

In two dimensions every refinement is a series move (the series/parallel
note, §2). A face of area \(A\) tiled by \(A/a^2\) plaquettes of weight
\(w_a\) has, after integrating the interior edges, the weight whose
coefficients are \(\hat c_R(a)^{A/a^2}\).

\textbf{Theorem 2.} Let \(w_a\) be symmetric class-function plaquette
weights on the lattice of spacing \(a\), with \(\hat c_R(a)\to1\) for
every \(R\). Suppose that for every \(R\) the limit

\[\psi(R)=\lim_{a\to0}\frac{1-\hat c_R(a)}{a^2}
=\lim_{a\to0}a^{-2}\int_G\Bigl(1-\frac{{\rm Re}\,\chi_R(g)}{d_R}\Bigr)
w_a(g)\,dg\]

exists. Then every face of area \(A\) has the limiting coefficients
\(e^{-A\psi(R)}\), and \(\psi\) has the form

\[\psi(R)=\frac{\sigma^2}2C_2(R)+\int_{G\setminus\{e\}}
\Bigl(1-\frac{{\rm Re}\,\chi_R(g)}{d_R}\Bigr)\nu(dg),\]

with \(\sigma^2\ge0\) and \(\nu\) a conjugation-invariant symmetric
measure satisfying

\[\int_G\min\bigl(1,d(e,g)^2\bigr)\,\nu(dg)<\infty.\] The continuum
theory is Yang--Mills with \(\lambda_2=\sigma^2\) exactly when
\(\nu=0\), and \(\nu=0\) holds when, for every \(\varepsilon>0\),
\(a^{-2}\int_{d(e,g)>\varepsilon}w_a\,dg\to0\).

\emph{Proof.} Since \(\hat c_R(a)\to1\),
\((A/a^2)\log\hat c_R(a)=-(A/a^2)(1-\hat c_R(a))(1+o(1))\to-A\psi(R)\).
Split the integral defining \(\psi\) at \(d(e,g)=\varepsilon\). Near the
identity, \(1-{\rm Re}\,\chi_R(e^X)/d_R=C_2(R)|X|^2/(2\dim G)+O(|X|^3)\)
after averaging over the directions of \(X\), which conjugation
invariance of \(w_a\) permits; this part converges to
\(\frac{\sigma^2}2C_2(R)\) with

\[\sigma^2=\lim_{\varepsilon\to0}\lim_{a\to0}\frac1{a^2\dim G}
\int_{d(e,g)<\varepsilon}|X|^2\,w_a(g)\,dg,\qquad g=e^X,\]

along a subsequence if needed. Away from the identity, the measures
\(a^{-2}w_a\,dg\) restricted to \(\{d\ge\varepsilon\}\) have bounded
mass (a finite sum of the functions \(1-{\rm Re}\,\chi_R/d_R\) over
representations separating the points of \(G\) is continuous and
positive off \(e\), hence bounded below there), so a subsequence
converges weakly to \(\nu\) on that set, conjugation invariant and
symmetric. The integrability of \(\nu\) near \(e\) follows from the
bound on the near part. The form of \(\psi\) agrees with Hunt's
classification of convolution semigroups on Lie groups, restricted to
central ones
(\href{https://doi.org/10.1090/S0002-9947-1956-0079232-9}{Hunt 1956},
metadata). The Lindeberg condition sets the far part to zero.
\(\square\)

The limiting objects are Lévy's two-dimensional Markovian holonomy
fields, which he constructs from Lévy processes on \(G\)
(\href{https://doi.org/10.24033/ast.785}{Lévy, Astérisque 329},
metadata). Theorem 2 adds only the lattice-side statement: which
plaquette actions lead to which member.

\textbf{Examples.} The Wilson action
\(w_a\propto e^{\beta_a{\rm Re\,tr}U}\) with
\(\beta_a\propto1/(\lambda_2a^2)\) concentrates like a Gaussian of width
\(\sqrt{\lambda_2}a\), satisfies the Lindeberg condition, and gives
Yang--Mills. A weight carrying an explicit centre-vortex fugacity,
\(w_a=(1-\kappa a^2)k_{\lambda_2a^2}+\frac{\kappa a^2}2\bigl[k_{\lambda_2a^2}(z\,\cdot) +k_{\lambda_2a^2}(z^{-1}\cdot)\bigr]\)
with \(z\) a generator of the centre \(\mathbb Z_N\) of \(SU(N)\), gives
\(\nu=\frac\kappa2(\delta_z+\delta_{z^{-1}})\) and

\[\psi(R)=\frac{\lambda_2}2C_2(R)+\kappa\Bigl(1-\cos\frac{2\pi k(R)}N\Bigr),\]

with \(k(R)\) the \(N\)-ality. On a spatial circle of circumference
\(L\) the refinement note's eq. (10) becomes
\(E_R-E_0=\hbar cL\,\psi(R)\): the vortex gas adds an
\(N\)-ality-dependent string tension. The limit is action-dependent
exactly through \((\sigma^2,\nu)\).

\hypertarget{any-dimension-the-error-budget-of-one-step}{%
\subsection{3. Any dimension: the error budget of one
step}\label{any-dimension-the-error-budget-of-one-step}}

Let \(t_n\) be the plaquette heat time at refinement step \(n\), which
by eq. (1) of the series/parallel note scales as \(\lambda_Da_n^{4-D}\)
at fixed physical coupling, and let \(L\) be a fixed physical box side.
One directional step leaves three kinds of error.

\begin{itemize}
\tightlist
\item
  \textbf{Perturbative.} Coupling shifts and irrelevant operators of
  relative size \(O(t_n)\) (Hypothesis P(\(\alpha\)) and Proposition 7
  of the series/parallel note).
\item
  \textbf{Large-field.} Configurations far from the identity at the
  lattice scale, with action at least \(c/t_n\), of density per physical
  volume about \(a_n^{-D}e^{-c/t_n}\); here \(c\) depends on the lattice
  action (for heat-kernel \(U(1)\) it is the vortex constant of Theorem
  5 there).
\item
  \textbf{Jump.} The action's Lévy content, the part of \(\nu\) in
  Theorem 2.
\end{itemize}

The continuum limit discards the lattice structure when the three are
summable over steps in a fixed physical box.

\begin{longtable}[]{@{}
  >{\raggedright\arraybackslash}p{(\columnwidth - 6\tabcolsep) * \real{0.2500}}
  >{\raggedright\arraybackslash}p{(\columnwidth - 6\tabcolsep) * \real{0.2500}}
  >{\raggedright\arraybackslash}p{(\columnwidth - 6\tabcolsep) * \real{0.2500}}
  >{\raggedright\arraybackslash}p{(\columnwidth - 6\tabcolsep) * \real{0.2500}}@{}}
\toprule\noalign{}
\begin{minipage}[b]{\linewidth}\raggedright
\(D\)
\end{minipage} & \begin{minipage}[b]{\linewidth}\raggedright
\(t_n\), fixed \(\lambda_D\)
\end{minipage} & \begin{minipage}[b]{\linewidth}\raggedright
Perturbative
\end{minipage} & \begin{minipage}[b]{\linewidth}\raggedright
Large-field, \(\sum(L/a_n)^De^{-c/t_n}\)
\end{minipage} \\
\midrule\noalign{}
\endhead
\bottomrule\noalign{}
\endlastfoot
2 & \(\lambda_2a_n^2\to0\) & summable & summable; limit set by
\((\sigma^2,\nu)\) \\
3 & \(\lambda_3a_n\to0\) & summable & summable faster than any power \\
4 & \(g^2(a_n)\to0\), logarithmically & resums into running &
\(\propto(a_n\Lambda)^{2b_0c-4}\), needs \(2b_0c>4\) \\
\(>4\) & \(\lambda_Da_n^{4-D}\to\infty\) & no small parameter & no
small-field regime \\
\end{longtable}

\textbf{The four-dimensional threshold.} With
\(g^{-2}(a)=2b_0\log(1/(a\Lambda))\), the convention of the refinement
note §6, \(e^{-c/g^2(a)}=(a\Lambda)^{2b_0c}\). For pure \(SU(N)\),
\(b_0=11N/(48\pi^2)\), so the large-field density per physical volume
vanishes iff

\[c>\frac2{b_0}=\frac{96\pi^2}{11N}:\qquad c>43.1\ (SU(2)),\qquad
c>28.7\ (SU(3)).\]

A continuum instanton has \(c=8\pi^2\approx79\), giving the density
\((a\Lambda)^{11N/3}a^{-4}\) of small instantons, which vanishes for
every \(N\ge2\) because \(11N/3>4\). Lattice actions admit dislocations,
small configurations of topological charge whose action lies below
\(8\pi^2\); whether they survive the continuum limit is decided by the
same inequality, the kind of criterion
\href{https://doi.org/10.1016/0370-2693(89)91067-8}{Pugh and Teper
(1989)} (metadata) analysed for \(SU(2)\), and the reason for
admissibility conditions on lattice fields
(\href{https://doi.org/10.1007/BF02029132}{Lüscher 1982}, metadata). The
numerical thresholds above are computed here from the one-loop
coefficient and are subject to the higher-order corrections of the
running.

\textbf{What this says about the mass gap.} In \(D=3\) at fixed
\(\lambda_3\) every plaquette action with vanishing jump content has an
ultraviolet limit governed by the same small-field analysis; a gap
\(C_3\hbar c\lambda_3\) is then a property of the infrared and, if it
exists, the same for all such actions. In \(D=4\) the same holds for
actions whose large-field constant exceeds \(2/b_0\), with gap
\(C_4\hbar c\Lambda\) and a \(\Lambda\) parameter that depends on the
action by a computable finite factor. The universality of the gap across
actions is thereby reduced to the three ultraviolet conditions, and the
gap itself remains the infrared obligation of the mass-gap map.

\hypertarget{consequence-for-state}{%
\subsection{4. Consequence for STATE}\label{consequence-for-state}}

Item 1 of STATE gains its action and dimension dependence. The next
steps on this line are to state Hypothesis P(\(\alpha\)) for a general
symmetric class-function action (second moment, large-field constant and
jump content as its data), and to connect the four-dimensional
large-field threshold to the H1 blocking hypothesis of the
\href{mass-gap-conditional-theorem.md}{conditional theorem}. For the
joint paper, Theorem 2 gives the two-dimensional row of the comparison
for every action at once, and the table gives the other rows.
\end{document}
